\chapter{Inference Fundamentals: From Logits to Tokens per Second}\label{ch:inference}

Everything we built in Parts~II--IV ends in a function that maps a sequence of token ids
to a matrix of logits. This chapter turns that function into a \emph{text generator}, and then
into a \emph{fast} text generator. The first half follows the path every practitioner walks the
first time they run an open-weight model: set up an environment, size the hardware, tokenize a
prompt, load a pretrained checkpoint (Qwen3-0.6B, the base model for all of Part~VI), and write
the autoregressive loop. The second half is where the systems person in you should lean in.
We will see that naive generation does a quadratic amount of redundant work, that the KV cache
removes it at the price of a memory footprint which quickly becomes \emph{the} resource that
limits a serving system, and that once the redundant work is gone, decoding is governed not by
FLOPs but by bytes: every generated token streams the entire model through the memory system.

The chapter closes by making the measurement vocabulary precise---time to first token, time per
output token, throughput, goodput---because every later chapter (batching engines, speculative
decoding, inference-time scaling, RL rollouts) is ultimately judged by these numbers.

\begin{figure}[t]
\centering
\begin{tikzpicture}[
  stage/.style={draw=sysblue,fill=sysblue!8,rounded corners=2pt,minimum width=2.3cm,
    minimum height=1.0cm,align=center,font=\small},
  sys/.style={stage,draw=napkin,fill=napkin!10},
  num/.style={circle,fill=sysblue,text=white,font=\scriptsize\bfseries,inner sep=1.2pt},
  arr/.style={-{Stealth[length=2mm]},thick,sysblue!70}]
\node[stage] (s1) at (0,0) {Environment\\\& hardware};
\node[stage] (s2) at (2.75,0) {Tokenize\\the prompt};
\node[stage] (s3) at (5.5,0) {Load\\Qwen3-0.6B};
\node[stage] (s4) at (8.25,0) {Autoregressive\\loop};
\node[stage] (s5) at (11.0,0) {Streaming\\generator};
\node[sys] (s6) at (11.0,-1.9) {KV cache};
\node[sys] (s7) at (8.25,-1.9) {Compilation\\\& CUDA graphs};
\node[sys] (s8) at (5.5,-1.9) {Prefill vs.\\decode roofline};
\node[sys] (s9) at (2.75,-1.9) {Latency\\metrics};
\node[sys,draw=labpurple,fill=labpurple!10] (s10) at (0,-1.9) {Lab 18};
\foreach \a/\b in {s1/s2,s2/s3,s3/s4,s4/s5,s5/s6,s6/s7,s7/s8,s8/s9,s9/s10} \draw[arr] (\a) -- (\b);
\foreach \n/\i in {s1/1,s2/2,s3/3,s4/4,s5/5,s6/6,s7/7,s8/8,s9/9,s10/10}
  \node[num] at (\n.north west) {\i};
\node[font=\scriptsize\itshape,text=sysblue] at (5.5,0.85) {make it work};
\node[font=\scriptsize\itshape,text=napkin] at (5.5,-2.75) {make it fast, and know \emph{why} it is fast};
\end{tikzpicture}
\caption{Roadmap of this chapter. The top row gets a pretrained model to produce text; the bottom row
removes redundant work and builds the performance model we use for the rest of the book.}
\label{fig:inference:roadmap}
\end{figure}

% =====================================================================
\section{From a prompt to a response}\label{sec:inference:pipeline}

A pretrained \emph{base} model has only seen next-token prediction on web-scale text. It has not
been instruction-tuned (\cref{ch:instruct}) or trained to reason (\cref{ch:grpo}), yet it already
produces fluent continuations and, surprisingly often, follows simple instructions, because
instruction--response pairs are common in its training data. Using such a model involves three
stages (\cref{fig:inference:pipeline}):

\begin{enumerate}
  \item \textbf{Input preparation.} The user's string is converted into a list of integer token
    ids by the tokenizer, optionally after wrapping it in a \emph{chat template}.
  \item \textbf{The model.} One or more forward passes map token ids to logits, i.e.\ unnormalized
    scores over the vocabulary for the next position.
  \item \textbf{Output processing.} A decoding rule (greedy argmax for now; temperature, top-$k$
    and top-$p$ in \cref{ch:scaling}) selects a token id, which is appended to the sequence; at the
    end, the tokenizer maps ids back to a string.
\end{enumerate}

\begin{figure}[ht]
\centering
\begin{tikzpicture}[scale=0.89,transform shape,
  box/.style={draw,rounded corners=2pt,minimum height=0.85cm,align=center,font=\small},
  tok/.style={draw=fdata,fill=fdata!12,minimum width=0.95cm,minimum height=0.55cm,font=\scriptsize\ttfamily},
  arr/.style={-{Stealth[length=2mm]},thick}]
% input lane
\node[font=\small\itshape,anchor=west] at (-0.2,0) {``Why is the sky blue?''};
\node[box,draw=fdata,fill=fdata!10,minimum width=2.4cm] (enc) at (5.6,0) {\code{tokenizer.encode}};
\foreach \t/\x in {Why/8.2,{\_is}/9.2,{\_the}/10.2,{\_sky}/11.2,{\_blue}/12.2,?/13.2}
  \node[tok] at (\x,0.0) {\t};
\draw[arr] (3.4,0) -- (enc.west);
\draw[arr] (enc.east) -- (7.65,0);
% model
\node[box,draw=fcompute,fill=fcompute!12,minimum width=6.3cm,minimum height=1.1cm] (llm) at (10.7,-1.7)
  {\textbf{LLM forward pass}\\[-1pt]{\scriptsize ids $(1,T)$ $\rightarrow$ logits $(1,T,V)$}};
\draw[arr] (10.7,-0.35) -- (llm.north);
% output lane
\node[box,draw=napkin,fill=napkin!12,minimum width=2.6cm] (sel) at (10.7,-3.4)
  {pick next id\\[-1pt]{\scriptsize argmax / sample}};
\draw[arr] (llm.south) -- (sel.north);
\node[box,draw=fdata,fill=fdata!10,minimum width=2.4cm] (dec) at (5.6,-3.4) {\code{tokenizer.decode}};
\draw[arr] (sel.west) -- (dec.east);
\node[font=\small\itshape,anchor=east] at (3.4,-3.4) {``\,The sky looks blue because\ldots''};
\draw[arr] (dec.west) -- (3.45,-3.4);
% feedback
\draw[arr,dashed,napkin] (sel.east) -- ++(1.9,0) |- (13.95,0) node[pos=0.25,right,font=\scriptsize,align=left]
  {append id,\\repeat};
\node[font=\scriptsize\bfseries,text=fdata,anchor=west] at (-0.2,0.6) {input preparation};
\node[font=\scriptsize\bfseries,text=fcompute,anchor=west] at (-0.2,-1.7) {model};
\node[font=\scriptsize\bfseries,text=napkin,anchor=west] at (-0.2,-2.8) {output processing};
\end{tikzpicture}
\caption{The three stages of text generation. Only the dashed feedback edge makes generation
\emph{autoregressive}; everything else is a pure function. (The underscore marks a leading space,
which most byte-level BPE tokenizers fold into the following token.)}
\label{fig:inference:pipeline}
\end{figure}

The model is the expensive part, but the other two stages are not free. In a serving system,
tokenization runs on the CPU and can become a bottleneck at high request rates; sampling over a
150k-entry vocabulary on the GPU involves sorts and softmaxes that are surprisingly costly when
fused poorly. We revisit both in \cref{ch:serving}.

\begin{keyidea}[Read diagrams bottom-up]
By convention, model diagrams are drawn with inputs at the bottom and outputs on top, so
information flows upward through the layer stack. We follow this convention for architecture
figures (\cref{fig:inference:qwen3stack}); data-flow figures such as \cref{fig:inference:pipeline}
read left to right.
\end{keyidea}

% =====================================================================
\section{Setting up the environment}\label{sec:inference:env}

All code for the book lives in the \ident{code/} directory: the importable package
\ident{llmsys} plus one directory per lab. Two setup paths work equally well.

\paragraph{pip.} With Python $\ge 3.10$ in any virtual environment:
\begin{lstlisting}[style=shell]
cd code
pip install -r requirements.txt     # torch, tokenizers, safetensors, sympy, ...
python -c "import llmsys, torch; print(torch.__version__)"
\end{lstlisting}
PyTorch wheels are hardware-specific. If the default wheel does not match your accelerator,
install PyTorch first using the selector on \texttt{pytorch.org}, then install the remaining
requirements.

\paragraph{uv.} The \texttt{uv} package manager resolves and installs dependencies much faster than pip,
creates an isolated environment automatically, and can even provision a Python interpreter.
After installing \texttt{uv}:
\begin{lstlisting}[style=shell]
cd code
uv venv && source .venv/bin/activate
uv pip install -r requirements.txt
uv run python -m pytest -q tests        # runs inside the managed environment
\end{lstlisting}
Either way, the environment must be \emph{reproducible}: pin versions in a lock file for
experiments you intend to report. Inference numerics depend on library versions (kernels,
reduction orders, TF32 defaults), and ``the same code'' can produce different greedy outputs
across PyTorch releases.

% =====================================================================
\section{Hardware: what you need and why}\label{sec:inference:hardware}

Training a frontier base model is out of reach for individuals. DeepSeek-V3, one of the few
recent models with a detailed compute disclosure, reports about 2.79 million H800 GPU-hours for its
full training on 14.8 trillion tokens, which the authors price at roughly US\$5.6 million at an
assumed rental rate of \$2 per GPU-hour, and that excludes all research and ablation runs.
This book therefore never pretrains a large model. We take a small but architecturally modern
pretrained model, Qwen3-0.6B, and apply to it the same post-training and inference-time
techniques that are used on models a hundred times larger. The methods transfer; only the
constants change.

Because chapters~\ref{ch:inference}--\ref{ch:refine} only run inference, they work on a laptop
CPU. Reinforcement learning (\cref{ch:grpo}) and distillation (\cref{ch:distill}) want a GPU.
The device-selection helper in \ident{llmsys.utils.get_device} picks CUDA, Apple's MPS, or the CPU,
in that order. On NVIDIA Ampere and newer GPUs it is also worth enabling TF32 tensor-core math
for float32 matmuls, which trades a few mantissa bits for a large speed-up:

\begin{lstlisting}[caption={Device selection and TF32.},label={lst:inference:device}]
import torch
from llmsys.utils import get_device

device = get_device()                 # cuda > mps > cpu
if device.type == "cuda":
    # fp32 matmuls run on tensor cores with a 10-bit mantissa (TF32)
    torch.backends.cuda.matmul.allow_tf32 = True
    torch.backends.cudnn.allow_tf32 = True
    print(torch.cuda.get_device_name(0))
print("using", device)
\end{lstlisting}

\begin{napkinmath}[Will it fit?]
A model with $N$ parameters stored in a format of $b$ bytes needs $Nb$ bytes for weights alone.
Qwen3-0.6B has $N=\num{596}$M parameters (the input embedding and the LM head are tied):
\begin{itemize}
  \item bf16: $0.596\times10^9 \times 2 \approx \SI{1.19}{GB}$; fp32: \SI{2.38}{GB}.
  \item Qwen3-8B: \SI{16.4}{GB} in bf16, which fits on a 24\,GB consumer card with a little room for the
    KV cache; Qwen3-32B (\SI{65.5}{GB}) needs an 80\,GB accelerator or quantization (\cref{ch:efficient}).
\end{itemize}
Weights are only the floor. The KV cache (\cref{sec:inference:kvcache}) adds memory proportional to
\emph{batch $\times$ context}, and for long reasoning traces it can exceed the weights.
\end{napkinmath}

Cloud notebooks (Google Colab, Lightning AI Studios and similar) are a convenient way to get a
GPU by the hour. For this chapter, a single entry-level data-center GPU such as an NVIDIA L4 is
more than enough; for RL in Part~VI an 80\,GB GPU makes life much easier.

\begin{systemslens}[Matching the device to the phase]
Different phases of the life cycle stress different resources. Pretraining and RL training are
dominated by dense matmuls (compute-bound, FLOP/s matter). Autoregressive decoding with small
batches is dominated by reading weights and KV cache (bandwidth-bound, TB/s matter). This is why
an RTX~4090 (\SI{165}{TFLOP/s} bf16, \SI{1.0}{TB/s}) can decode a small model at a respectable
fraction of an H100's speed (\SI{989}{TFLOP/s}, \SI{3.35}{TB/s}) even though it has one sixth the
compute: for batch-1 decoding the relevant ratio is bandwidth, roughly 1:3.3. Apple-silicon
laptops with large unified memory and high bandwidth are popular local-inference machines for the
same reason.
\end{systemslens}

% =====================================================================
\section{Preparing inputs: tokenizers and chat templates}\label{sec:inference:tokenizer}

Every prompt, every reasoning trace and every answer must become a list of integers before
the model can see it. Qwen3 uses a byte-level BPE tokenizer (the algorithm we implemented in
\cref{ch:text}) with a vocabulary of $V=151{,}936$ entries. We load it from the
\code{tokenizer.json} file shipped with the checkpoint using the Hugging Face \ident{tokenizers}
library, wrapped by \ident{llmsys.qwen3.Qwen3Tokenizer}.

\begin{lstlisting}[caption={A tokenizer round trip.},label={lst:inference:roundtrip}]
from llmsys.qwen3 import Qwen3Tokenizer

tok = Qwen3Tokenizer("qwen3-0.6b-base/tokenizer.json")
prompt = "Why is the sky blue?"
ids = tok.encode(prompt)            # list[int]
assert tok.decode(ids) == prompt     # lossless for byte-level BPE
for i in ids:
    print(f"{i:>7d} -> {tok.decode([i])!r}")
\end{lstlisting}

Printing ids one at a time shows the subword structure: frequent words are single tokens (usually
including their leading space), rarer words are split into pieces, and \emph{nothing} is ever out of
vocabulary, because any byte sequence can fall back to single-byte tokens. Made-up words and
non-English text simply tokenize into more pieces (\cref{ex:inference:unknown}).

\paragraph{Vocabulary size is a systems trade-off.} GPT-2 used about 50k tokens, Llama~3 about
128k, Qwen3 about 152k. A larger vocabulary compresses text into fewer tokens, so every request
is shorter (fewer decode steps, smaller KV cache), but it makes the embedding table and the output
projection larger.

\begin{napkinmath}[The cost of the LM head]
For Qwen3-0.6B, the output projection is a $1024\times151{,}936$ matrix: \SI{155.6}{M}
parameters, 26\% of the whole model. Per generated token it costs
$2\cdot1024\cdot151{,}936\approx\SI{311}{MFLOP}$ out of $\approx\SI{1.19}{GFLOP}$ total. During
decoding it is also \SI{311}{MB} of bf16 weights that must be read for every step. Small models
pay disproportionately for big vocabularies; for Qwen3-8B the head is only about 8\% of
parameters. Two practical consequences: (i)~during prefill, compute the LM head only for the last
position (our serving engine in \cref{ch:serving} does this); (ii)~sampling kernels must handle
$V$-length rows efficiently.
\end{napkinmath}

\paragraph{Chat templates.} Instruction-tuned and reasoning models are trained on conversations
serialized with special tokens. Qwen uses the ChatML layout:
\begin{lstlisting}[style=plain]
<|im_start|>user
Why is the sky blue?<|im_end|>
<|im_start|>assistant
\end{lstlisting}
The model is expected to continue after the final \code{assistant} header and to end its turn
with \code{<|im\_end|>}. The \ident{apply_chat_template} method of our tokenizer wrapper produces this string.
For the reasoning variant of Qwen3 the template controls ``thinking mode'': an empty
\code{<think>}\ldots\code{</think>} block pre-filled into the assistant turn disables the reasoning trace.
A base model has no special relationship with these tokens, so in \cref{ch:eval} we will use
a plain-text prompt format for it instead. Getting the template wrong is the single most common
cause of mysteriously bad outputs from open-weight models.

\begin{pitfall}[Template and EOS mismatches]
Base Qwen3 ends documents with \code{<|endoftext|>} (id 151643), while chat-tuned variants end
turns with \code{<|im\_end|>}. If you stop on the wrong id, a base model happily continues past the
end of its answer into a new, unrelated document. Always check which token(s) your checkpoint
uses to terminate, and pass all of them to the stopping logic.
\end{pitfall}

% =====================================================================
\section{Loading pretrained weights}\label{sec:inference:loading}

We use Qwen3-0.6B for three reasons: it is small enough for consumer hardware, it is architecturally
identical to its much larger siblings (\cref{ch:modern}), and it ships both as a \emph{base}
model (our starting point for reasoning training) and as an official \emph{reasoning} variant
that serves as a reference in evaluations. Our implementation in \ident{llmsys.qwen3} is plain
PyTorch, independent of any LLM library, yet weight-compatible with the official checkpoint.

\begin{lstlisting}[caption={Loading Qwen3-0.6B from a downloaded snapshot.},label={lst:inference:load}]
# one-time download (needs internet):
#   huggingface-cli download Qwen/Qwen3-0.6B-Base --local-dir qwen3-0.6b-base
import torch
from llmsys.qwen3 import load_qwen3_from_dir

model = load_qwen3_from_dir("qwen3-0.6b-base", size="0.6B",
                            dtype=torch.bfloat16, device=device)
print(model)       # 28 x Qwen3Block(GQA + SwiGLU + 2 RMSNorm)
\end{lstlisting}

\ident{load_qwen3_from_dir} constructs \ident{Qwen3Model} with randomly initialized placeholders,
reads every \code{.safetensors} shard, and copies tensors in by name (the mapping is in
\ident{load_qwen3_hf_state_dict}). Printing the module shows the structure summarized in
\cref{fig:inference:qwen3stack}: a $151{,}936\times1024$ embedding, 28 identical blocks, and a final
norm and output head. Each block's attention has 16 query heads but only 8 key/value heads of
dimension 128, so \code{W\_query} maps $1024\to2048$ while \code{W\_key}/\code{W\_value} map $1024\to1024$.
The MLP is a SwiGLU with hidden size 3072. We derived all of these pieces in \cref{ch:modern}; here
the model is a black box that turns ids into logits.

\begin{figure}[ht]
\centering
\begin{tikzpicture}[
  blk/.style={draw,rounded corners=2pt,minimum width=4.6cm,minimum height=0.55cm,font=\small,align=center},
  arr/.style={-{Stealth[length=1.6mm]},thick},
  note/.style={font=\scriptsize,align=left,text=black!70}]
\node[blk,fill=fdata!12] (emb) at (0,0) {Token embedding $151936\times1024$};
\node[blk,fill=keygreen!8] (n1) at (0,1.05) {RMSNorm};
\node[blk,fill=fcompute!15,minimum height=0.8cm] (att) at (0,2.05) {Grouped-query attention\\[-2pt]{\scriptsize 16 Q heads, 8 KV heads, RoPE, QK-norm}};
\node[circle,draw,inner sep=0.5pt] (p1) at (0,3.0) {\scriptsize $+$};
\node[blk,fill=keygreen!8] (n2) at (0,3.75) {RMSNorm};
\node[blk,fill=fcompute!15,minimum height=0.8cm] (ff) at (0,4.75) {SwiGLU MLP\\[-2pt]{\scriptsize $1024\to3072\to1024$}};
\node[circle,draw,inner sep=0.5pt] (p2) at (0,5.7) {\scriptsize $+$};
\node[blk,fill=keygreen!8] (fn) at (0,6.6) {Final RMSNorm};
\node[blk,fill=fdata!12] (head) at (0,7.45) {LM head $1024\times151936$ (tied)};
\node[font=\small] (out) at (0,8.3) {logits $(B,T,151936)$};
\node[font=\small] (in) at (0,-0.85) {token ids $(B,T)$};
\draw[arr] (in) -- (emb); \draw[arr] (emb) -- (n1); \draw[arr] (n1) -- (att); \draw[arr] (att) -- (p1);
\draw[arr] (p1) -- (n2); \draw[arr] (n2) -- (ff); \draw[arr] (ff) -- (p2); \draw[arr] (p2) -- (fn);
\draw[arr] (fn) -- (head); \draw[arr] (head) -- (out);
\draw[arr,black!60] (0,0.5) -- (-2.75,0.5) |- (p1.west);
\draw[arr,black!60] (0,3.3) -- (-2.75,3.3) |- (p2.west);
\draw[decorate,decoration={brace,amplitude=5pt},thick,sysblue] (2.55,6.05) -- (2.55,0.55)
  node[pos=0.25,right=7pt,note,text=sysblue]{\textbf{$\times$28} identical\\blocks};
\node[note,anchor=west] at (2.5,7.45) {155.6M params (26\%)\\311 MFLOP / token};
\node[note,anchor=west] at (3.3,2.05) {KV cache lives here:\\$2\times8\times128$ values\\per token per layer};
\end{tikzpicture}
\caption{Qwen3-0.6B as a data-flow stack (bottom to top). Grey arrows are residual connections.
The annotations point at the two places that matter most for inference cost: the large tied
output head and the per-layer KV state.}
\label{fig:inference:qwen3stack}
\end{figure}

\begin{systemslens}[Loading is an I/O problem]
Loading a checkpoint is a copy from disk to host memory to device memory. A \SI{1.19}{GB} bf16
checkpoint takes about \SI{0.4}{s} at \SI{3}{GB/s} NVMe speed and about \SI{0.05}{s} over a
PCIe~4.0 x16 link (\SI{\approx25}{GB/s} effective); for a 70B model the same numbers become tens of
seconds, which matters for autoscaling and for RL systems that reload weights frequently
(\cref{ch:rlsys}). Three tricks are standard:
\begin{itemize}
  \item \textbf{safetensors + mmap:} a flat header plus raw little-endian tensor bytes; tensors can
    be memory-mapped without unpickling and without an intermediate copy (unlike \code{torch.load}
    of a pickle, which can also execute arbitrary code).
  \item \textbf{Meta-device construction:} build the module on \code{torch.device("meta")} so no
    memory is allocated for random initial weights, then materialize directly from the checkpoint
    (\code{load\_state\_dict(..., assign=True)}).
  \item \textbf{Direct-to-GPU and sharded loading:} each tensor-parallel rank reads only its slice
    (\cref{ch:efficient}); GPUDirect Storage can bypass host memory entirely.
\end{itemize}
\end{systemslens}

% =====================================================================
\section{The autoregressive loop, one step at a time}\label{sec:inference:loop}

LLMs generate one token at a time. At each iteration the model sees the prompt plus everything
generated so far, predicts one new token, and the token is appended to the input of the next
iteration (\cref{fig:inference:unrolled}). This is \emph{autoregressive} generation.

\begin{figure}[ht]
\centering
\begin{tikzpicture}[
  t/.style={draw=black!40,fill=white,minimum width=0.62cm,minimum height=0.48cm,font=\scriptsize,inner sep=0pt},
  p/.style={t,fill=fdata!15},
  g/.style={t,fill=napkin!25},
  n/.style={t,fill=napkin!70,text=white,font=\scriptsize\bfseries},
  llm/.style={draw=fcompute,fill=fcompute!12,rounded corners=2pt,minimum height=0.6cm,font=\scriptsize\bfseries},
  arr/.style={-{Stealth[length=1.5mm]},thick}]
\foreach \r/\k in {0/0,1/1,2/2,3/3} {
  \pgfmathsetmacro{\y}{-1.35*\r}
  \node[font=\scriptsize,anchor=east] at (-0.4,\y) {iter.\ \the\numexpr\r+1\relax};
  \foreach \i in {0,...,4} \node[p] at (0.65*\i,\y) {};
  \ifnum\k>0 \foreach \j in {1,...,\k} \node[g] at (0.65*4+0.65*\j,\y) {}; \fi
  \node[llm,minimum width=1.3cm] (L\r) at (8.0,\y) {LLM};
  \node[n] (N\r) at (9.6,\y) {};
  \draw[arr] (0.65*4+0.65*\k+0.45,\y) -- (L\r.west);
  \draw[arr] (L\r) -- (N\r);
}
\foreach \r in {0,1,2} {
  \pgfmathsetmacro{\y}{-1.35*\r}
  \pgfmathsetmacro{\yy}{-1.35*(\r+1)}
  \pgfmathsetmacro{\xx}{0.65*4+0.65*(\r+1)}
  \draw[arr,napkin,dashed] (N\r.south) to[out=-90,in=90,looseness=0.6] (\xx,\yy+0.25);
}
\node[font=\scriptsize,text=fdata] at (1.3,0.55) {prompt (5 tokens)};
\node[font=\scriptsize,text=napkin] at (9.6,0.55) {new token};
\node[font=\scriptsize,align=left,anchor=west,text=black!70] at (10.1,-2.0)
  {the input grows\\by one token per\\iteration};
\end{tikzpicture}
\caption{Autoregressive generation unrolled over four iterations. Each output token (dark) is fed
back as the last input token of the next iteration (light). Without further tricks, iteration
$i$ re-processes all $5+i-1$ tokens.}
\label{fig:inference:unrolled}
\end{figure}

Zoom into one iteration. A causal LM processes the whole input in a single forward pass and
returns one vector of $V$ logits \emph{per input position}: the logits at position $t$ score the
candidates for position $t+1$. For a $T$-token input the output has shape $(1, T, V)$. During
training all $T$ predictions are used (that is what makes pretraining efficient, \cref{ch:pretrain}).
During generation, the predictions for positions inside the prompt are about tokens we already
have, so only the \emph{last} row matters (\cref{fig:inference:lastrow}).

\begin{figure}[ht]
\centering
\begin{tikzpicture}[cell/.style={draw=black!25,minimum width=0.55cm,minimum height=0.42cm,inner sep=0pt,font=\tiny}]
\foreach \r in {0,...,5} {
  \foreach \c in {0,...,9} {
    \pgfmathsetmacro{\shade}{int(mod(37*\r+17*\c,60))}
    \ifnum\r=5 \node[cell,fill=napkin!\shade!white] at (0.55*\c,-0.42*\r) {}; \else
    \node[cell,fill=black!\shade!white,opacity=0.35] at (0.55*\c,-0.42*\r) {}; \fi
  }
}
\node[cell,fill=napkin,draw=napkin!60!black,line width=1pt] at (0.55*6,-0.42*5) {};
\node[font=\scriptsize,anchor=east] at (-0.45,-0.42*2.5) {$T=6$ positions};
\node[font=\scriptsize] at (0.55*4.5,0.45) {$V=151{,}936$ vocabulary columns (10 shown)};
\draw[decorate,decoration={brace,amplitude=4pt,mirror},thick] (5.45,-0.42*4.5) -- (5.45,0.2)
  node[midway,right=6pt,font=\scriptsize,align=left]{predictions for tokens\\already in the prompt:\\discarded at inference};
\draw[-{Stealth},thick,napkin] (5.45,-0.42*5) -- (7.0,-0.42*5) node[right,font=\scriptsize,align=left]
  {\code{logits[:, -1]}\\$\to$ \code{argmax} $\to$ next id};
\draw[-{Stealth},thick,napkin] (0.55*6,-0.42*5-0.3) -- (0.55*6,-2.85) node[below,font=\scriptsize]{largest score: the greedy choice};
\end{tikzpicture}
\caption{One forward pass returns a $T\times V$ score matrix (batch dimension dropped). Generation
reads only the last row and selects the column with the largest score.}
\label{fig:inference:lastrow}
\end{figure}

\begin{lstlisting}[caption={One step of generation by hand.},label={lst:inference:onestep}]
ids = torch.tensor(tok.encode(prompt), device=device).unsqueeze(0)   # (1, T)
with torch.inference_mode():
    logits = model(ids)                 # (1, T, V)
print(logits.shape, logits.dtype)       # torch.Size([1, T, 151936]) torch.bfloat16
last = logits[0, -1]                    # (V,) scores for position T+1
next_id = torch.argmax(last, dim=-1, keepdim=True)   # (1,)
print(repr(tok.decode(next_id.tolist())))
\end{lstlisting}

A few PyTorch details in \cref{lst:inference:onestep} are worth spelling out:
\begin{itemize}
  \item \code{unsqueeze(0)} adds the batch dimension the model expects; \code{squeeze(0)} removes a
    size-1 dimension again. Shape bookkeeping is the main source of bugs in generation code.
  \item \code{torch.max} returns the largest \emph{value}, \code{torch.argmax} its \emph{index}. The index
    into the vocabulary dimension \emph{is} the token id. \code{keepdim=True} keeps a trailing
    dimension of size 1, which is convenient for concatenating the id onto a $(B,T)$ tensor.
  \item Always choosing the argmax is \emph{greedy decoding}: deterministic, cheap, and usually
    fine for short factual answers, but prone to repetition. \Cref{ch:scaling} introduces sampling.
  \item \code{torch.inference\_mode()} is a stricter, faster version of \code{torch.no\_grad()}: it
    disables autograd tracking \emph{and} version counters on tensors, so nothing is saved for a
    backward pass. The logits come back in bf16 because the weights are bf16.
\end{itemize}

\begin{deepdive}[``Inference'' is an overloaded word]
In statistics, \emph{inference} means estimating unknown quantities from data, with uncertainty.
In deep learning it means something narrower: evaluating a trained, frozen function on new inputs.
When we say ``inference'' in this book we mean the latter. Nothing is learned while generating;
the model is a fixed map from token ids to next-token scores, and the decoding rule is a fixed map
from scores to an id.
\end{deepdive}

% =====================================================================
\section{A minimal streaming generator}\label{sec:inference:generator}

Wrapping the single step in a loop gives a complete generator. We write it as a Python
\emph{generator function} that \code{yield}s every token as soon as it exists. Streaming is not
cosmetic: users perceive latency as ``time until text starts appearing'', and every serving API
(server-sent events, websockets) exposes tokens incrementally.

\begin{lstlisting}[caption={Greedy streaming generation without a cache (Lab 18, Task 1).},label={lst:inference:nocache}]
@torch.inference_mode()
def generate_nocache(model, token_ids, max_new_tokens, eos_id=None):
    model.eval()                                   # disable dropout etc.
    for _ in range(max_new_tokens):
        logits = model(token_ids)[:, -1]           # (B, V): last position only
        nxt = torch.argmax(logits, dim=-1, keepdim=True)   # (B, 1)
        if eos_id is not None and bool((nxt == eos_id).all()):
            return                                 # every row has finished
        yield nxt
        token_ids = torch.cat([token_ids, nxt], dim=1)     # grow the input
\end{lstlisting}

The caller decodes and prints each token as it arrives:
\begin{lstlisting}
ids = torch.tensor([tok.encode("Explain KV caching in one sentence.")], device=device)
for t in generate_nocache(model, ids, max_new_tokens=100, eos_id=tok.pad_token_id):
    print(tok.decode(t[0].tolist()), end="", flush=True)    # flush: show it now
\end{lstlisting}

Two observations about the output of a base model. First, the first emitted token usually starts
with a space, because BPE attaches leading spaces to words; strip it if you care. Second, without
a stopping condition a base model finishes its answer, emits the document separator
\code{<|endoftext|>}, and continues into an unrelated document. Passing the separator's id as
\code{eos\_id} stops generation at the right place. In a batch, we can only stop when \emph{every} row
has finished, which is the first hint of the batching inefficiencies we fix in \cref{ch:serving}.

\paragraph{Measuring it.} We need a stopwatch that is honest on asynchronous hardware. CUDA
kernel launches return immediately, so a timer that does not call \code{torch.cuda.synchronize()}
measures launch overhead, not execution. \Cref{lst:inference:measure} records a timestamp after
each yielded token and derives the metrics we define precisely in \cref{sec:inference:metrics};
peak GPU memory is available from \code{torch.cuda.max\_memory\_allocated()}.

\begin{lstlisting}[caption={Timing a streaming generator (Lab 18, Task 6).},label={lst:inference:measure}]
def measure_generation(gen_fn, model, token_ids, max_new_tokens,
                       clock=time.perf_counter):
    _sync(); t0 = clock(); stamps = []
    for _ in gen_fn(model, token_ids, max_new_tokens):
        _sync()                     # wait for the GPU before reading the clock
        stamps.append(clock())
    n = len(stamps)
    ttft = stamps[0] - t0                               # includes the prefill
    tpot = (stamps[-1] - stamps[0]) / (n - 1)           # mean inter-token gap
    return dict(n_tokens=n, ttft=ttft, tpot=tpot,
                tokens_per_s=n / (stamps[-1] - t0))
\end{lstlisting}


\paragraph{Streaming detokenization.} Decoding tokens one at a time is subtly wrong. A byte-level
BPE token is a sequence of \emph{bytes}, not characters: an emoji or a CJK character is 3--4 UTF-8
bytes and is frequently split across two or three tokens. Calling \code{decode([t])} on each piece
produces replacement characters (\texttt{U+FFFD}) for the incomplete fragments. Production servers
decode \emph{incrementally}: keep the ids generated so far, decode a small trailing window, and only
emit the part of the text that can no longer change (no trailing replacement character); the rest
is held back until the next token arrives.

\begin{lstlisting}[caption={Incremental detokenization that never emits partial characters.},label={lst:inference:detok}]
class IncrementalDetokenizer:
    def __init__(self, tok, window=8):
        self.tok, self.ids, self.emitted, self.window = tok, [], 0, window

    def push(self, token_id) -> str:
        self.ids.append(token_id)
        start = max(0, len(self.ids) - self.window)   # bounded work per token
        prefix = self.tok.decode(self.ids[start:-1]) if len(self.ids) > 1 else ""
        full = self.tok.decode(self.ids[start:])
        if full.endswith("\ufffd"):                  # incomplete UTF-8 sequence
            return ""                                # wait for more bytes
        new = full[len(prefix):] if full.startswith(prefix) else full
        return new
\end{lstlisting}
The window keeps the cost per token constant (decoding the full history every step would be
quadratic, the same mistake as recomputing without a KV cache), and a serving front end typically
runs this on the CPU, overlapped with the next GPU decode step.

On a laptop CPU, \cref{lst:inference:nocache} with Qwen3-0.6B produces a handful of tokens per second,
and it gets slower as the answer gets longer. The rest of the chapter explains why, and fixes it.

% =====================================================================
\section{KV caching}\label{sec:inference:kvcache}

Look again at \cref{fig:inference:unrolled}. In iteration $i$ the model recomputes the hidden
states of all earlier tokens, although they are bit-for-bit identical to those of iteration $i-1$:
in a causal model, position $t$ never attends to positions after $t$, so adding a token at the end
cannot change anything that came before. The only thing a new query needs from earlier positions
are their \emph{keys} and \emph{values} in every attention layer. Store them, and each decode step
only has to process the single new token.

\begin{figure}[ht]
\centering
\begin{tikzpicture}[
  t/.style={draw=black!40,minimum width=0.5cm,minimum height=0.42cm,inner sep=0pt,font=\tiny},
  arr/.style={-{Stealth[length=1.5mm]},thick}]
% left: no cache
\node[font=\small\bfseries] at (2.0,1.0) {without cache: recompute};
\foreach \r in {0,...,5} {
  \pgfmathtruncatemacro{\n}{3+\r}
  \foreach \c in {1,...,\n} {
    \ifnum\c=\n \node[t,fill=fcompute!70] at (0.5*\c-0.5,-0.5*\r) {};
    \else \node[t,fill=warnred!25] at (0.5*\c-0.5,-0.5*\r) {}; \fi
  }
  \node[font=\tiny,anchor=east] at (-0.4,-0.5*\r) {step \the\numexpr\r+1\relax};
}
\node[font=\scriptsize,align=center] at (2.0,-3.4) {work per step $\propto$ prefix length\\total $\approx \tfrac{1}{2}n^2$ token-passes};
% right: cache
\begin{scope}[xshift=7.0cm]
\node[font=\small\bfseries] at (2.0,1.0) {with KV cache: reuse};
\foreach \r in {0,...,5} {
  \pgfmathtruncatemacro{\n}{3+\r}
  \foreach \c in {1,...,\n} {
    \ifnum\c=\n \node[t,fill=fcompute!70] at (0.5*\c-0.5,-0.5*\r) {};
    \else \node[t,fill=fmemory!25] at (0.5*\c-0.5,-0.5*\r) {}; \fi
  }
}
\node[font=\scriptsize,align=center] at (2.0,-3.4) {compute per step: 1 token;\\reads all cached K,V (memory $\propto$ length)};
\end{scope}
\node[t,fill=warnred!25] at (0.6,-4.3) {}; \node[font=\scriptsize,anchor=west] at (0.9,-4.3) {recomputed (wasted)};
\node[t,fill=fmemory!25] at (4.6,-4.3) {}; \node[font=\scriptsize,anchor=west] at (4.9,-4.3) {read from cache};
\node[t,fill=fcompute!70] at (8.2,-4.3) {}; \node[font=\scriptsize,anchor=west] at (8.5,-4.3) {newly computed};
\end{tikzpicture}
\caption{Work per decode step. Without a cache, the redundant (red) triangle dominates.
With a cache, compute per step is constant, but the cached state that must be \emph{read} grows
linearly, which is why long-context decoding becomes memory-bound.}
\label{fig:inference:kvwork}
\end{figure}

Our attention modules already support this: they accept a \code{cache} object and, if present,
call \code{cache.update(layer\_idx, k\_new, v\_new)}, which appends the new keys and values and returns
the full history (\cref{ch:attention,ch:modern}). The generator changes in exactly two places:

\begin{lstlisting}[caption={Greedy generation with a KV cache (Lab 18, Task 2).},label={lst:inference:cache}]
from llmsys.kvcache import KVCache

@torch.inference_mode()
def generate_kvcache(model, token_ids, max_new_tokens, eos_id=None):
    model.eval()
    cache = KVCache(model.n_layers)
    logits = model(token_ids, cache=cache)[:, -1]       # PREFILL: whole prompt once
    for _ in range(max_new_tokens):
        nxt = torch.argmax(logits, dim=-1, keepdim=True)
        if eos_id is not None and bool((nxt == eos_id).all()):
            return
        yield nxt
        logits = model(nxt, cache=cache)[:, -1]         # DECODE: one token per step
\end{lstlisting}

The first call, with the whole prompt, is called \textbf{prefill}; each subsequent single-token
call is a \textbf{decode} step. Positional information must stay consistent: Qwen3 rotates queries
and keys by their \emph{absolute} position (RoPE), so the attention module asks the cache how many
tokens it already holds and uses that as the position offset; GPT-2 offsets its learned position
embeddings the same way. The critical sanity check is that cached and uncached generation produce
\emph{identical} greedy tokens; Lab~18 tests exactly that.

\paragraph{Complexity.} Generating $n$ tokens after a $P$-token prompt without a cache costs
$\sum_{t=0}^{n-1}(P+t)\approx nP+\tfrac{1}{2}n^2$ token-passes through the dense layers. With a
cache it costs $P+n$. For $P=32$, $n=512$ that is a factor of $\approx270$ fewer dense-layer FLOPs.
Attention itself is still linear per step (the new query attends over $P+t$ cached keys), so the
total attention work is quadratic in both cases, but with a tiny constant.

\begin{napkinmath}[How big is the KV cache?]
Per token, every layer stores one key and one value vector per KV head:
\begin{equation}\label{eq:inference:kvbytes}
  \text{bytes/token} = 2 \cdot L \cdot H_{kv} \cdot d_h \cdot b .
\end{equation}
\begin{itemize}
  \item Qwen3-0.6B ($L{=}28$, $H_{kv}{=}8$, $d_h{=}128$, bf16): $2\cdot28\cdot8\cdot128\cdot2=\SI{114688}{B}=\SI{112}{KiB}$
    per token. A single 32k-token reasoning trace: \SI{3.8}{GB}, three times the weights. The full
    40,960-token context: \SI{4.7}{GB}.
  \item Had Qwen3-0.6B used full multi-head attention (16 KV heads), it would need \SI{224}{KiB}/token:
    grouped-query attention halves the cache (\cref{ch:modern}).
  \item Qwen3-8B: \SI{144}{KiB}/token; Qwen3-32B: \SI{256}{KiB}/token. A batch of 64 sequences of 8k
    tokens on Qwen3-8B needs \SI{77}{GB} of KV cache, about five times its \SI{16.4}{GB} of weights.
\end{itemize}
The KV cache, not the weights, is what bounds batch size and context length in a serving system.
\end{napkinmath}

\begin{systemslens}[Growing vs.\ static caches]
Our \ident{KVCache} grows by \code{torch.cat}, which allocates a new tensor and copies the entire
history at every step: $O(t)$ bytes moved at step $t$, $O(n^2)$ in total, plus allocator churn.
\ident{StaticKVCache} preallocates $(B,H_{kv},T_{\max},d_h)$ buffers once and writes each new
slice in place: no reallocation, no copies, and fixed memory addresses. Our version still returns a
growing \emph{view} \code{[:, :, :end]}, so attention sees a different key length every step; fully
shape-static engines attend over the whole $T_{\max}$ buffer and mask positions beyond the current
length. Fixed shapes and addresses are what make CUDA graphs and compiled kernels reusable (next section). They come at a price:
memory is reserved for the maximum length up front, even if the request stops early.
\Cref{ch:serving} resolves this tension with paged allocation.
\end{systemslens}

% =====================================================================
\section{Compiling the model}\label{sec:inference:compile}

Eager PyTorch executes one Python operator at a time: each \code{nn.Linear}, each RMSNorm, each
elementwise multiply dispatches through the Python interpreter, the dispatcher, and finally launches
a kernel. For a large batch this overhead is invisible. For batch-1 decoding of a small model, the
kernels are tiny and the CPU-side overhead can dominate: a Qwen3-0.6B decode step launches several
hundred kernels, and at a few microseconds each the launch path alone can rival the
$\approx$\SI{0.4}{ms} it takes an H100 to read the weights.

\code{torch.compile} attacks this from two directions (\cref{fig:inference:compile}).
\emph{TorchDynamo} traces Python bytecode into an FX graph guarded by assumptions (shapes,
dtypes, Python constants); \emph{Inductor} lowers the graph, fuses chains of elementwise and
reduction ops into single generated (Triton or C++) kernels, and picks matmul implementations.
With \code{mode="reduce-overhead"} the result is additionally captured into a \emph{CUDA graph}, so a
whole decode step is replayed with a single launch.

\begin{figure}[ht]
\centering
\begin{tikzpicture}[
  st/.style={draw,rounded corners=2pt,minimum height=0.95cm,align=center,font=\scriptsize,minimum width=2.6cm},
  arr/.style={-{Stealth[length=1.6mm]},thick}]
\node[st,fill=black!5] (py) at (0,0) {Python model\\\code{forward()}};
\node[st,fill=fcompute!12] (dyn) at (3.0,0) {TorchDynamo\\bytecode $\to$ FX graph\\+ guards};
\node[st,fill=fcompute!12] (aot) at (6.0,0) {AOTAutograd\\functionalize\\(+ backward)};
\node[st,fill=fcompute!12] (ind) at (9.0,0) {Inductor\\fusion, codegen,\\matmul choice};
\node[st,fill=keygreen!12] (krn) at (12.1,0.7) {fused Triton /\\C++ kernels};
\node[st,fill=keygreen!12] (cg) at (12.1,-0.8) {CUDA graph\\(reduce-overhead)};
\draw[arr] (py) -- (dyn); \draw[arr] (dyn) -- (aot); \draw[arr] (aot) -- (ind);
\draw[arr] (ind.east) -- (krn.west); \draw[arr] (ind.east) -- (cg.west);
\draw[arr,warnred,dashed] (dyn.south) -- ++(0,-0.75) node[below,font=\scriptsize,text=warnred,align=center]
  {guard fails (new shape)\\$\Rightarrow$ recompile};
\draw[arr,napkin,dashed] (dyn.north) -- ++(0,0.55) node[above,font=\scriptsize,text=napkin]{graph break $\Rightarrow$ fall back to eager};
\end{tikzpicture}
\caption{The \code{torch.compile} stack. Static shapes avoid recompilation and enable CUDA graph
capture; data-dependent Python control flow causes graph breaks.}
\label{fig:inference:compile}
\end{figure}

\begin{lstlisting}[caption={Compiling and benchmarking with warm-up runs.},label={lst:inference:compile}]
compiled = torch.compile(model)           # or mode="max-autotune" / "reduce-overhead"
for rep in range(3):                      # rep 0 pays for tracing + codegen
    m = measure_generation(generate_static, compiled, ids, 100)
    tag = "warm-up" if rep == 0 else f"timed run {rep}"
    print(f"{tag}: {m['tokens_per_s']:.1f} tok/s, TTFT {m['ttft']*1e3:.0f} ms")
\end{lstlisting}

Three practical rules follow from \cref{fig:inference:compile}:
\begin{enumerate}
  \item \textbf{Always warm up.} The first call traces and generates code and can take tens of
    seconds. Report only subsequent runs, and say so.
  \item \textbf{Keep shapes static on the hot path.} A growing \code{torch.cat} cache gives every
    decode step a different key length; Dynamo then either recompiles or falls back to dynamic
    shapes, which limits fusion and rules out CUDA graphs. A static buffer attended with a length
    mask plus a position counter held in a device tensor gives identical shapes every step. Modules that mutate Python integers
    (\code{self.pos += 1}) can also trigger recompiles; recent PyTorch versions offer configuration
    flags to treat such integers as dynamic rather than baked-in constants.
  \item \textbf{Prefill and decode are different programs.} Prefill has shape $(B,P)$ with variable
    $P$; decode has shape $(B,1)$. Engines compile (and graph-capture) decode for a fixed set of
    batch sizes and pad the live batch up to the nearest captured size, while prefill runs eagerly
    or with dynamic shapes.
\end{enumerate}

The speed-up from compilation depends heavily on where the time goes. On a CPU, fusion reduces
memory traffic and Python overhead; with a cache the per-step work is so small that overhead
matters a lot, and compiled cached decoding is typically several times faster than eager. Without
a cache the work is dominated by large matmuls that compilation cannot make faster, so the gain is
small. If default compilation gives nothing, try \code{mode="max-autotune"}, which benchmarks
alternative kernels and tile sizes at compile time.

\begin{pitfall}[Benchmarking the compiler instead of the model]
Comparing a compiled model's \emph{first} run against an eager model's steady state, forgetting
\code{synchronize()}, letting the GPU clock ramp up during the first measurement, or measuring
with a 10-token output where TTFT dominates are the four most common ways to publish meaningless
inference numbers. Lab~18 makes you control all four.
\end{pitfall}


% =====================================================================
\section{Where the time goes: launch overhead and CUDA graphs}\label{sec:inference:overhead}

The roofline bound of the next section is a statement about the GPU. Eager PyTorch at batch 1 often
never gets close to it, because the GPU is waiting for the CPU. Count the kernels in one Qwen3
decode step:

\begin{napkinmath}[Kernels per decode step]
Per block, an eager forward pass launches roughly: two RMSNorms ($\approx$6 small kernels each:
cast, square, mean, rsqrt, two multiplies), three QKV GEMMs, two QK-norms ($\approx$12), RoPE on
$q$ and $k$ (slicing, \code{cat}, two multiplies and an add each, $\approx$10), the cache update
(2), \code{repeat\_interleave} for GQA (2), scores GEMM, scale, mask, softmax, value GEMM ($\approx$6),
output projection, residual add, three SwiGLU GEMMs plus SiLU and multiply, and a second residual
add. That is $\approx$50 kernels per block, about 1,400 for 28 blocks, plus the head.
At 4--8\,\si{\micro s} of CPU time per eager op (Python, dispatcher, launch), the CPU needs
\SIrange{6}{11}{ms} to \emph{issue} a step whose GPU work takes \SI{0.4}{ms} at the bandwidth bound.
Batch-1 eager decoding on a fast GPU is therefore \emph{launch-bound}: the GPU idles about 95\% of the time.
\end{napkinmath}

Three remedies, in increasing order of effectiveness: (i) \emph{fusion} (\code{torch.compile},
hand-written fused RMSNorm/RoPE kernels, \cref{app:triton}) reduces the number of kernels;
(ii) \emph{asynchrony}: as long as the CPU runs ahead, launches overlap with GPU work, so avoid
anything that forces a sync on the hot path (\code{.item()}, \code{.tolist()}, Python \code{if} on a
GPU tensor, as in our EOS check!); (iii) \emph{CUDA graphs} record the entire sequence of kernel
launches of one step once, with fixed memory addresses, and replay it with a single
\code{cudaGraphLaunch}. Graphs require static shapes and static buffers, which is exactly what
\ident{StaticKVCache} provides.

\begin{lstlisting}[caption={Capturing one static-cache decode step in a CUDA graph (GPU only).},label={lst:inference:cudagraph}]
# Assumes: model on cuda, cache = StaticKVCache(...) already holding the prefill,
# and a model/cache that read the write position from a *device tensor*, so the
# replayed graph does not bake in a Python integer (an exercise in Lab 18).
static_in = torch.zeros(B, 1, dtype=torch.long, device="cuda")
s = torch.cuda.Stream(); s.wait_stream(torch.cuda.current_stream())
with torch.cuda.stream(s):                 # warm up on a side stream
    for _ in range(3):
        static_out = model(static_in, cache=cache)
torch.cuda.current_stream().wait_stream(s)
g = torch.cuda.CUDAGraph()
with torch.cuda.graph(g):                  # record kernels + addresses, no execution
    static_out = model(static_in, cache=cache)
# decode loop: copy the new token into the static input, replay, read static output
for _ in range(n):
    static_in.copy_(next_ids)
    g.replay()                             # ~1 launch instead of ~1,400
    next_ids = static_out[:, -1].argmax(-1, keepdim=True)
\end{lstlisting}

The comment in \cref{lst:inference:cudagraph} hides the real engineering: our cache tracks its
length in a Python list, and Python state is frozen into a captured graph. Production engines keep
positions, sequence lengths and slot indices in device tensors that are updated by \code{copy\_}
before each replay, and capture one graph per supported batch size (e.g.\ 1, 2, 4, 8, \ldots, 256),
padding the live batch up to the next captured size. \code{torch.compile(mode="reduce-overhead")}
automates exactly this bookkeeping for you.

\begin{systemslens}[Overhead hides in the sampler too]
Our EOS check \code{bool((nxt == eos\_id).all())} copies a GPU boolean to the host and blocks
until the step finishes, which serializes CPU and GPU every token. Engines avoid per-token host
syncs by keeping stop conditions on the device, checking them only every $k$ steps, or running
the scheduler for step $t+1$ while the GPU executes step $t$ (``async scheduling''). On a modern GPU
these few microseconds per token matter: at \SI{0.4}{ms} per step, a \SI{50}{\micro s} sync is a
12\% slowdown.
\end{systemslens}

% =====================================================================
\section{Anatomy of a decode step: the roofline view}\label{sec:inference:roofline}

We now have the vocabulary to answer the question this chapter promised: what limits the speed of
generation? Recall the roofline model from \cref{ch:pytorch}: a kernel that performs $F$ FLOPs and
moves $M$ bytes takes at least $\max(F/\Pi,\ M/\beta)$, where $\Pi$ is peak compute and $\beta$
peak memory bandwidth. Its \emph{arithmetic intensity} $I=F/M$ decides which term wins; the
crossover (the \emph{ridge point}) is $\Pi/\beta$, about 295 FLOP/byte on an H100 (989 TFLOP/s over
3.35 TB/s), 153 on an A100 and 164 on an RTX~4090.

\paragraph{Weight GEMMs.} A linear layer with an $m\times k$ weight applied to $B$ token vectors
does $2Bmk$ FLOPs and must read $mk\,b$ bytes of weights (activations are comparatively tiny when
$B\ll k$). Its intensity is
\begin{equation}\label{eq:inference:intensity}
  I_{\text{GEMM}} \approx \frac{2Bmk}{mk\,b} = \frac{2B}{b} \quad\text{FLOP/byte},
\end{equation}
i.e.\ exactly $B$ for bf16. A decode step processes one token per sequence, so $B$ is the
\emph{batch size}; a prefill step processes $P$ tokens per sequence, so $B$ is the \emph{number of
prompt tokens}. Prefill of a 2,000-token prompt has intensity $\approx2000$: firmly compute-bound.
Batch-1 decode has intensity $\approx1$: two orders of magnitude below the ridge, so the tensor cores
sit idle more than 99\% of the time while HBM streams weights.

\paragraph{Attention over the cache.} A decode query attends over $C$ cached tokens: $\approx4\,C\,H\,d_h$
FLOPs per layer to read $2\,C\,H_{kv}\,d_h\,b$ bytes. Its intensity is $\approx(H/H_{kv})\cdot(2/b)$,
i.e.\ the GQA group size for bf16 (2 for Qwen3-0.6B), and \emph{independent of batch size}: each
sequence reads its own cache. Batching amortizes weight reads but never KV reads.

\begin{figure}[ht]
\centering
\begin{tikzpicture}
\begin{loglogaxis}[width=0.92\linewidth,height=7.0cm,xmin=0.5,xmax=20000,ymin=1e12,ymax=2e15,
  xlabel={arithmetic intensity (FLOP/byte)},ylabel={attainable FLOP/s},
  grid=major,grid style={black!10},legend columns=2,
  legend style={font=\scriptsize,at={(0.5,-0.2)},anchor=north,draw=black!30},
  tick label style={font=\scriptsize},label style={font=\small}]
\addplot[very thick,sysblue,domain=0.5:20000,samples=200] {min(989e12, 3.35e12*x)};
\addlegendentry{H100 roofline (bf16)}
\addplot[thick,dashed,napkin,domain=0.5:20000,samples=200] {min(165e12, 1.008e12*x)};
\addlegendentry{RTX 4090 roofline}
\addplot[only marks,mark=*,mark size=2.2pt,warnred] coordinates {(1,3.35e12) (8,26.8e12) (64,214.4e12) (256,857.6e12)};
\addlegendentry{decode, weight GEMMs ($B$=1,8,64,256)}
\addplot[only marks,mark=square*,mark size=2.4pt,keygreen] coordinates {(2048,989e12) (512,989e12)};
\addlegendentry{prefill (512 / 2048 tokens)}
\addplot[only marks,mark=diamond*,mark size=3pt,labpurple] coordinates {(2,6.7e12)};
\addlegendentry{decode attention (GQA, group 2)}
\draw[dotted,thick] (axis cs:295,1e12) -- (axis cs:295,2e15);
\node[font=\scriptsize,rotate=90,anchor=south] at (axis cs:295,2e13) {ridge $\approx$ 295};
\end{loglogaxis}
\end{tikzpicture}
\caption{Roofline placement of inference work on an H100 (analytic, from published peak numbers).
Decode GEMMs sit at intensity $\approx B$ and climb the bandwidth slope as the batch grows; attention
over the KV cache stays near the GQA group size regardless of batch; prefill is compute-bound.}
\label{fig:inference:roofline}
\end{figure}

\begin{napkinmath}[A lower bound on decode speed]
A decode step must read all weights once plus every sequence's KV cache. With
$N$ parameters, $b$ bytes each, batch $B$, context $C$, and $\kappa$ KV bytes per token:
\[
  t_{\text{step}} \;\ge\; \max\!\left(\frac{2NB}{\Pi},\ \frac{Nb + B\,C\,\kappa}{\beta}\right).
\]
For Qwen3-0.6B on an H100 at $B=1$, $C=1024$: $(1.19\,\text{GB}+0.12\,\text{GB})/3.35\,\text{TB/s}\approx\SI{0.39}{ms}$,
i.e.\ at most $\approx2{,}560$ tokens/s for a single stream, and compute would need only
\SI{1.2}{\micro s}. For Qwen3-8B the bound is \SI{4.9}{ms} ($\approx200$ tokens/s); for Qwen3-32B
\SI{19.6}{ms} ($\approx50$ tokens/s). Real systems reach a good fraction of these bounds only with
CUDA graphs and fused kernels; eager PyTorch at batch 1 is often launch-bound and far slower.
\end{napkinmath}

\begin{table}[ht]
\centering\small
\caption{Roofline lower bounds for one decode step on an H100 at context 1,024 (bf16, attention FLOPs ignored).
Throughput rises almost linearly with $B$ while weight reads dominate, then saturates as per-sequence KV
reads take over. Computed with \ident{llmsys.flops.decode_step_time_lower_bound}.}
\label{tab:inference:decode}
\begin{tabular}{rrrrrrr}
\toprule
 & \multicolumn{3}{c}{Qwen3-0.6B} & \multicolumn{3}{c}{Qwen3-8B}\\
\cmidrule(lr){2-4}\cmidrule(lr){5-7}
$B$ & step (ms) & tokens/s & bound & step (ms) & tokens/s & bound\\
\midrule
1   & 0.39 & 2{,}560   & memory & 4.94  & 203    & memory\\
8   & 0.64 & 12{,}600  & memory & 5.25  & 1{,}520  & memory\\
64  & 2.60 & 24{,}600  & memory & 7.77  & 8{,}230  & memory\\
256 & 9.33 & 27{,}400  & memory & 16.4  & 15{,}600 & memory\\
512 & 18.3 & 28{,}000  & memory & 28.0  & 18{,}300 & memory\\
\bottomrule
\end{tabular}
\end{table}

\Cref{tab:inference:decode} contains the central lesson of LLM serving. For the small model, KV
reads overtake weight reads already at $B\approx10$ (each sequence's 1,024-token cache is
\SI{115}{MB}, a tenth of the weights), so batching saturates early; for the 8B model weights dominate
for longer and batching pays off over a wider range. Either way, throughput is bought with batch
size, and batch size is bought with KV-cache memory. Everything in \cref{ch:serving,ch:efficient}
(paged KV memory, prefix sharing, quantized caches, GQA/MLA, speculative decoding) is a way to move
along this trade-off curve.

\begin{systemslens}[Prefill and decode want different machines]
Prefill is compute-bound and its latency scales with prompt length; decode is bandwidth-bound and
its latency scales with model size and context. Running both on the same GPU means a long prefill
stalls every decoding user for its duration, while decode steps leave the tensor cores idle.
Production systems either interleave the phases carefully (chunked prefill), or
\emph{disaggregate} them onto separate GPU pools and ship the KV cache between them. Both are
designed in \cref{ch:serving}.
\end{systemslens}

% =====================================================================
\section{Measuring inference: metrics and honest benchmarks}\label{sec:inference:metrics}

Throughput alone does not describe a serving system, because users experience latency, and the two
pull in opposite directions. The standard vocabulary (\cref{fig:inference:timeline}):

\begin{description}
  \item[TTFT] (time to first token): from request arrival to the first output token. It contains
    queueing delay plus the prefill, and grows with prompt length. Interactive applications target
    a few hundred milliseconds.
  \item[TPOT / ITL] (time per output token / inter-token latency): the gap between consecutive
    tokens during decoding. Reading speed is around 5--10 tokens/s, so a TPOT below
    \SI{100}{ms} feels instantaneous for chat; agents and reasoning models that emit thousands of
    hidden tokens need far lower TPOT.
  \item[E2E latency:] $\text{TTFT} + (n-1)\cdot\text{TPOT}$ for $n$ output tokens.
  \item[Throughput:] output tokens per second across all users (sometimes input+output tokens),
    the quantity that determines cost per token.
  \item[Goodput:] the throughput of requests that \emph{meet their latency SLOs}, e.g.\ ``TTFT
    $<$ \SI{500}{ms} and TPOT $<$ \SI{50}{ms} at P90''. A system can raise throughput by batching
    harder while its goodput collapses.
\end{description}

\begin{figure}[ht]
\centering
\begin{tikzpicture}[x=1cm,y=1cm]
\draw[-{Stealth},thick] (0,0) -- (13,0) node[right,font=\scriptsize]{time};
\fill[black!10] (0.3,0.15) rectangle (2.2,0.75); \node[font=\scriptsize] at (1.25,0.45) {queue};
\fill[fcompute!35] (2.2,0.15) rectangle (4.6,0.75); \node[font=\scriptsize] at (3.4,0.45) {prefill};
\foreach \x in {4.6,5.7,6.8,7.9,9.0,10.1,11.2} {
  \fill[napkin!60] (\x,0.15) rectangle (\x+0.12,0.75);
  \draw[napkin!70!black] (\x+0.12,0.15) -- (\x+0.12,0.75);
}
\foreach \x/\l in {4.72/1,5.82/2,6.92/3,8.02/4,9.12/5,10.22/6,11.32/7}
  \node[font=\tiny,above] at (\x,0.75) {$y_{\l}$};
\draw[|-|] (0.3,-0.35) -- (4.72,-0.35) node[midway,below,font=\scriptsize]{TTFT};
\draw[|-|] (4.72,-0.35) -- (5.82,-0.35) node[midway,below,font=\scriptsize]{TPOT};
\draw[|-|] (0.3,-1.15) -- (11.32,-1.15) node[midway,below,font=\scriptsize]{end-to-end latency};
\node[font=\scriptsize,anchor=west] at (0.3,1.3) {request arrives};
\draw[-{Stealth}] (0.3,1.15) -- (0.3,0.8);
\node[font=\scriptsize,anchor=east] at (11.32,1.3) {last token / EOS};
\draw[-{Stealth}] (11.32,1.15) -- (11.32,0.8);
\end{tikzpicture}
\caption{Latency anatomy of one streamed request. Queueing and prefill determine TTFT; the decode
cadence determines TPOT. Under load, other users' prefills interleaved with our decode steps make
TPOT jittery, which is why tail percentiles matter.}
\label{fig:inference:timeline}
\end{figure}

\paragraph{Little's law.} In steady state, the number of requests in the system equals arrival
rate times mean latency: $\bar{N} = \lambda\,\bar{W}$. If a server sustains a decode batch of 64
and each request takes \SI{10}{s} end to end, it can admit at most 6.4 requests/s; pushing the arrival
rate beyond that only grows the queue and therefore TTFT. This simple identity is the right first
model for capacity planning, and we use it again with the simulator in \cref{ch:serving}.

\begin{keyidea}[A benchmarking checklist]
\begin{enumerate}
  \item Synchronize the device before every clock read; use \code{time.perf\_counter}.
  \item Warm up (compilation, autotuning, allocator, GPU clocks) and discard warm-up runs.
  \item Report TTFT and TPOT separately, with the prompt and output lengths used.
  \item Report distributions (P50/P90/P99), not just means, and the number of repetitions.
  \item Fix the output length (ignore EOS or force it) when comparing systems, otherwise a model
    that stops earlier looks faster.
  \item State hardware, dtype, batch size, library versions and whether results are compiled.
\end{enumerate}
\end{keyidea}

\paragraph{What we measured.} \Cref{fig:inference:measured} shows per-step decode time for a small
Qwen3-shaped model (8 layers, $d=512$, random weights) running eagerly on the two-core CPU of the
machine this chapter was written on. Absolute numbers are meaningless for a GPU, but the
\emph{shapes} are exactly the ones predicted above: without a cache, step time grows linearly with
sequence length (so total time grows quadratically); with either cache, step time is nearly flat
and dominated by fixed per-step costs at these lengths. You will reproduce this plot, and the
batch-size and compilation sweeps, on your own hardware in Lab~18.

\begin{figure}[ht]
\centering
\begin{tikzpicture}
\begin{groupplot}[group style={group size=2 by 1,horizontal sep=1.6cm},width=0.49\linewidth,height=5.6cm,
  grid=major,grid style={black!10},tick label style={font=\scriptsize},label style={font=\small},
  legend style={font=\scriptsize,at={(0.03,0.97)},anchor=north west}]
\nextgroupplot[xlabel={sequence length (tokens)},ylabel={ms per decode step (median)},ymin=0,
  title={\small per-step latency vs.\ length}]
\addplot[thick,mark=*,mark size=1.3pt,warnred] coordinates {(48,128.06) (80,168.89) (112,197.91) (144,239.8) (176,215.92) (208,226.78) (240,307.75) (272,347.13) (304,552.38) (336,647.84) (368,473.57) (400,414.49) (432,431.66) (464,1047.08) (496,841.88)};
\addlegendentry{no cache}
\addplot[thick,mark=square*,mark size=1.3pt,sysblue] coordinates {(48,26.28) (80,30.8) (112,26.5) (144,31.6) (176,26.74) (208,27.59) (240,26.12) (272,28.07) (304,34.52) (336,29.61) (368,35.35) (400,32.81) (432,31.68) (464,35.81) (496,34.08)};
\addlegendentry{growing KV cache}
\addplot[thick,mark=triangle*,mark size=1.6pt,keygreen] coordinates {(48,28.62) (80,28.12) (112,29.45) (144,29.44) (176,33.94) (208,26.17) (240,24.6) (272,22.58) (304,32.24) (336,28.64) (368,30.87) (400,31.72) (432,33.41) (464,35.05) (496,30.25)};
\addlegendentry{static KV cache}
\nextgroupplot[xmode=log,log basis x=2,xlabel={batch size $B$},ylabel={aggregate decode tokens/s},ymin=0,
  title={\small static-cache decoding vs.\ batch},xtick={1,2,4,8,16,32},xticklabels={1,2,4,8,16,32}]
\addplot[thick,mark=*,mark size=1.3pt,napkin] coordinates {(1,42.8) (2,81.4) (4,116.3) (8,152.8) (16,204.3) (32,377.8)};
\addlegendentry{measured}
\addplot[dashed,black!50,domain=1:32] {42.8*x};
\addlegendentry{perfect scaling}
\end{groupplot}
\end{tikzpicture}
\caption{Measured, not modelled: a 6-layer, $d{=}384$ Qwen3-shaped model with random weights, fp32,
eager PyTorch, one thread on a heavily shared two-core CPU (hence the noise). Left: 480 tokens after a
32-token prompt; total time was \SI{203}{s} without a cache and \SI{15}{s} with either cache. Right:
aggregate throughput grows with batch size, sub-linearly because even on a CPU larger batches
eventually stop being bound by reading the weights.}
\label{fig:inference:measured}
\end{figure}

% =====================================================================
\section{Summary}\label{sec:inference:summary}
\begin{itemize}
  \item Text generation is a loop: tokenize, run the model, keep only the last position's logits,
    select a token, append, repeat; stream tokens as they are produced and stop on the right EOS id.
  \item The tokenizer and chat template are part of the model's interface; vocabulary size trades
    shorter sequences against a larger embedding and LM head.
  \item The KV cache removes the quadratic recomputation of naive decoding; its size is
    $2LH_{kv}d_hb$ bytes per token (\SI{112}{KiB} for Qwen3-0.6B) and it, not the weights, limits batch
    size and context length.
  \item Prefill is compute-bound; decode is memory-bound with arithmetic intensity $\approx$ batch size
    for the weight GEMMs and $\approx$ GQA group size for attention.
  \item \code{torch.compile} removes Python/launch overhead and fuses kernels; static shapes (static
    cache) enable CUDA graphs; always warm up.
  \item Eager batch-1 decoding is launch-bound on fast GPUs ($\sim$1,400 kernels per step); fusion,
    avoiding host syncs, and CUDA graphs close the gap to the bandwidth bound.
  \item Measure TTFT, TPOT, throughput and goodput separately, with synchronization and percentiles.
\end{itemize}

% =====================================================================
\section{Exercises}\label{sec:inference:exercises}

\begin{exercise}[Unknown words]\label{ex:inference:unknown}
Encode five invented words and two sentences in a language other than English with the Qwen3
tokenizer. How many tokens per character do you get compared with ordinary English text? Why can
a byte-level BPE tokenizer never produce an ``unknown'' token, and what is the cost of rare
scripts in terms of KV-cache memory per character?
\end{exercise}

\begin{exercise}[Running on an accelerator]
Run Lab~18's \code{run.py} on a GPU (or Apple silicon) and on the CPU. Report TTFT, TPOT and
tokens/s for the three decoders, eager and compiled. Which variant gains the most from moving to
the GPU, and why?
\end{exercise}

\begin{exercise}[Counting redundant work]
For a prompt of $P$ tokens and $n$ generated tokens, derive the exact number of token-positions
processed by the dense layers with and without a KV cache, and the number of query--key dot
products in attention in both cases. Evaluate the ratios for $P=1000$, $n=100$ and for $P=100$,
$n=1000$. Which workload benefits more from caching?
\end{exercise}

\begin{exercise}[KV budget planning]
An 80\,GB H100 serves Qwen3-8B in bf16. Reserve \SI{4}{GB} for activations and the CUDA context.
(a) How many tokens fit in the KV cache? (b) How many concurrent 8k-token conversations? (c) How
does the answer change with an fp8 KV cache, and with int4 weights (\cref{ch:efficient})?
\end{exercise}

\begin{exercise}[Critical batch size]
Using \cref{eq:inference:intensity}, compute the batch size at which the weight GEMMs of a decode
step become compute-bound on an A100, H100 and RTX~4090 for bf16 weights and for int8 weights with
bf16 compute. Why does weight-only quantization \emph{lower} the critical batch size?
\end{exercise}

\begin{exercise}[The LM head at prefill]
Our model computes logits for all $T$ positions. For a 4,096-token prefill on Qwen3-0.6B, how many
FLOPs and how much activation memory (bf16) does the full logits tensor cost, versus computing only
the last row? Modify a copy of \ident{Qwen3Model.forward} to accept a \code{last\_only} flag and
measure the difference.
\end{exercise}

\begin{exercise}[Cache copies]
Instrument \ident{KVCache.update} to count bytes copied by \code{torch.cat}. Show that the total over
$n$ decode steps is quadratic in $n$, and compare with \ident{StaticKVCache}. At what sequence length
does the copy traffic exceed the weight traffic for Qwen3-0.6B?
\end{exercise}

\begin{exercise}[Recompilation hunting]
Compile the model and run \code{generate\_kvcache} (growing cache) with
\code{TORCH\_LOGS=recompiles}. How many recompilations happen in 50 steps? Repeat with
\code{generate\_static}. Explain the difference in terms of Dynamo guards.
\end{exercise}

\begin{exercise}[Little's law]
A service has mean output length 300 tokens, TPOT \SI{25}{ms}, TTFT \SI{300}{ms}, and a maximum decode batch of
128. Estimate the maximum sustainable request rate. What happens to TTFT if requests arrive 20\%
faster than that rate for ten minutes?
\end{exercise}

\begin{exercise}[Determinism]
Generate 200 tokens greedily with batch size 1 and again as row 0 of a batch of 8 (other rows
random prompts of the same length). Are the outputs identical? Explain why batch-dependent kernel
selection and reduction order can change greedy outputs, and why this matters for evaluation
reproducibility (\cref{ch:eval}) and for RL (\cref{ch:rlsys}).
\end{exercise}

% =====================================================================
\section{Lab 18: Three decoders and a roofline}\label{sec:inference:lab}
\begin{labbox}{18}{Three decoders and a roofline}
\labdir{lab18_inference}
\labsection{Goals}
Implement cache-less, dynamically cached and statically cached greedy decoding; verify they produce
identical tokens; quantify KV memory and the decode roofline; and measure TTFT/TPOT correctly.

\labsection{Setup}
\begin{lstlisting}[style=shell]
cd code
python3 -m pytest -q labs/lab18_inference          # tests your starter.py
\end{lstlisting}
Tests use tiny GPT-2 and Qwen3 configurations with random weights and run on a CPU in seconds.
\code{run.py} uses the real Qwen3-0.6B base checkpoint (download command in its header).

\labsection{Tasks}
\begin{enumerate}
  \item \ident{generate_nocache}: streaming greedy decoding that reprocesses the full sequence.
  \item \ident{generate_kvcache}: one prefill call, then single-token decode calls with \ident{KVCache}.
  \item \ident{generate_static}: the same with a preallocated \ident{StaticKVCache} of length $T+n$.
  \item \ident{kv_bytes_per_token} and \ident{max_tokens_in_budget} for both GPT-2 (MHA) and Qwen3 (GQA) configs.
  \item \ident{decode_roofline} and \ident{critical_batch} implementing the napkin model of \cref{sec:inference:roofline}.
  \item \ident{measure_generation}: TTFT, TPOT, total time and tokens/s with device synchronization
    (the tests inject a fake clock).
\end{enumerate}

\labsection{Measurements}
\begin{itemize}
  \item \code{bench.py --exp length --plot len.png}: per-step latency vs.\ sequence length for the
    three decoders. Expect linear growth without a cache and a flat line with either cache.
  \item \code{bench.py --exp batch}: TPOT and aggregate tokens/s vs.\ batch size. Find the knee and
    compare it with \ident{critical_batch} for your device (on a GPU with \code{--size 0.6B}).
  \item \code{bench.py --exp compile}: eager vs.\ \code{torch.compile} default vs.\
    \code{reduce-overhead}, three repetitions each. Expect the first repetition to be slow and the
    cached, static-shape decoder to benefit most.
  \item \code{run.py}: the real model; confirm all three decoders emit identical text.
\end{itemize}

\labsection{Deliverables}
Passing tests; three plots (length, batch, compile) with axes labelled and hardware stated; a
half-page note comparing your measured batch-1 decode speed with the roofline bound and explaining
the gap.

\labsection{Stretch goals}
\begin{itemize}
  \item Capture the static-cache decode step in a CUDA graph manually (\code{torch.cuda.CUDAGraph})
    and compare against \code{reduce-overhead}.
  \item Add a \code{last\_only} flag to the model forward and measure prefill memory savings.
  \item Plot measured tokens/s against \ident{decode_roofline} predictions over batch sizes on your GPU
    and report the achieved fraction of bandwidth.
\end{itemize}
\end{labbox}
